\subsection{From Imitation Error to Task Performance}
\label{sec:actionanalysis}
An action deviation during exposure can leave the target inaccessible or displace it before execution.
We first relate coordinated action error to complete-task failure, then transfer imitation risk from demonstrations to deployment through a learned representation.

\subsubsection{Task-Performance Difference Between Policies}
\label{sec:actionconsequences}
We compare a policy $\pi$ with a reference policy $\pi_0$ using a finite-horizon performance-difference identity~\cite{ross2014aggrevate,cheng2018value}.
At each history, the identity measures the consequence of selecting an alternative action and then following $\pi_0$.
Given a history $\mathcal H_t=h$ and a candidate action
$a\in\mathcal A$, define the action-value function as the
probability of failing the complete exposure-and-execution task ($\Psi=0$):
\begin{equation}
 Q_t^{\pi_0}(h,a)
 :=\Prob^{a,\pi_0}\!\left(\Psi=0\mid\mathcal H_t=h\right).
 \label{eq:historyactionvalue}
\end{equation}
Here, $\Prob^{a,\pi_0}$ takes action $a$ at time $t$ and follows $\pi_0$ afterward, so an exposure action is evaluated through its consequences for the remaining task.

When the current action is also selected by $\pi_0$, this reduces
to the value function of $\pi_0$, $V_t^{\pi_0}(h) :=Q_t^{\pi_0}\bigl(h,\pi_0(h)\bigr).$
The effect of choosing $a$ instead of $\pi_0(h)$ is therefore
measured by
\begin{equation}
 A_t^{\pi_0}(h,a)
 :=Q_t^{\pi_0}(h,a)-V_t^{\pi_0}(h).
 \label{eq:actionconsequence}
\end{equation}
Thus, $A_t^{\pi_0}(h,a)>0$ indicates an increase in failure probability, whereas $A_t^{\pi_0}(h,a)<0$ indicates a decrease.

The comparison requires a common task and environment.
\begin{assumption}[Common task and environment]
\label{ass:commonprocess}
The compared policies share the same initial distribution and
task-completion criterion. Given a history $h$ and an action $a$,
the next available history follows a conditional distribution
that does not otherwise depend on the policy.
The reference policy $\pi_0$ is defined on all histories
involved in the comparison.
\end{assumption} 
% 比较的是同一任务实例分布与相同交互观测机制，而非强制相同后续轨迹。
The policies may induce different trajectories, but the initial task distribution and interaction and sensing mechanisms remain fixed.

Let $F(\pi)=1-J(\pi)$ denote the task-failure probability of policy $\pi$.
To evaluate $\pi$ relative to $\pi_0$, we set $a=\pi(h)$
and aggregate the resulting action consequences along trajectories generated by $\pi$.

Applying the finite-horizon performance-difference identity~\cite{kakade2002,ross2014aggrevate,cheng2018value}
to these one-step action effects gives
\begin{equation}
 F(\pi)-F(\pi_0)
 =\E_\pi\!\left[
 \sum_{t=0}^{T-1}A_t^{\pi_0}(\mathcal H_t,\pi(\mathcal H_t))
 \right].
\label{eq:pd}
\end{equation}
The expectation is over trajectories generated by $\pi$, so the action effects are evaluated at histories visited by the evaluated policy.
Appendix~\ref{app:pdidentity} gives the proof.


\subsubsection{Task-Failure Bounds from Imitation Risk}
\label{sec:actionerrorbound}
Taking the expert policy as the reference,
we show that the imitation risk $R_\nu(\pi)$ bounds task-failure probability $F(\pi)$
when demonstrations sufficiently cover exposure and execution histories encountered at deployment
and action deviations have bounded task consequences.

Imitation risk is evaluated under the demonstration distribution $\nu$,
but deployment histories may follow a different distribution because actions depend on $\pi$.
Let $\mu_\pi^t$ denote the distribution of $(t,\mathcal H_t)$ at control step $t$ under $\pi$.
The distribution $\bar\mu_\pi=T^{-1}\sum_{t=0}^{T-1}\mu_\pi^t$ corresponds to choosing
a control step uniformly and observing its history.

% 对待评价策略，假设部署分布由示范分布以有限常数控制。
\begin{assumption}[Demonstration coverage]
\label{ass:democoverage}

For the evaluated policy $\pi$, there exists a uniform constant
$1\le C<\infty$ such that
\[
\bar\mu_\pi(\mathcal B)\le C\,\nu(\mathcal B)
\]
for every measurable set $\mathcal B$ of pairs $(t,h)$.
\end{assumption}
This excludes deployment histories absent from the demonstration distribution and limits probability underestimation after occluder and target interactions.

\begin{assumption}[Action sensitivity]
\label{ass:actionsensitivity}
Let $A_t^{\mathrm E}$ denote the action effect with the expert
policy as the reference.
For a fixed action scale, there exists a uniform constant
$0\le L<\infty$ such that
\[
|A_t^{\mathrm E}(h,a)|
\le L\norm{a-\pi_{\mathrm E}(h)}
\]
for all relevant $(t,h,a)$.
\end{assumption}
This translates coordinated command error into a change in eventual failure probability.
Deformable interaction does not guarantee such regularity near changes in contact or accessibility.

Following the connection between deployment imitation error
and task cost~\cite{ross2011}, Assumptions~\ref{ass:democoverage}
and~\ref{ass:actionsensitivity}, together with the
performance-difference identity, yield the following bound (Appendix~\ref{app:actionriskbound}):
\begin{equation}
 F(\pi)\le F(\pi_{\mathrm E})+TL\sqrt{C R_\nu(\pi)}.
 \label{eq:successbound}
\end{equation}

Here, $C$ measures demonstration-to-deployment shift and $L$ measures task sensitivity to action deviations.
The bound shows that imitation accuracy controls task failure only under both conditions.

\subsubsection{Imitation-Risk Transfer Through Representations}
\label{sec:representationcoverage}
We next replace complete-history coverage with coverage in a learned representation and account for the error introduced by merging histories.

Assumption~\ref{ass:democoverage} can fail even after an arbitrarily small action deviation.
For $T\ge2$, suppose the expert always executes $u_0=0$, whereas the learned policy executes
an admissible $u_0=\varepsilon\ne0$.
Histories recording $u_0=\varepsilon$ then occur under the learned policy
but have zero probability under the expert demonstration distribution.
No finite coverage constant $C$ can satisfy the assumption.

A representation can merge such distinctions, potentially relaxing coverage requirements,
but may also discard information needed to distinguish action consequences.
For exposure and execution, useful compression should suppress irrelevant appearance variation without merging situations that require further exposure with those that permit target interaction.
This tension between distributional alignment and information loss under a noninvertible representation
also appears in domain-adaptation analyses~\cite{johansson2019support,zhao2019invariant}.

Fix the policy and expert, set $\mu=\bar\mu_\pi$, and let $Z=\phi(t,\mathcal H_t)$ be a measurable causal representation that retains the time index.
Let $\mu_Z$ and $\nu_Z$ be the representation distributions induced by $\mu$ and $\nu$.

\begin{assumption}[Representation coverage]
\label{ass:representationcoverage}
There exists $1\le C_\phi<\infty$ such that
\begin{equation}
 \mu_Z(\mathcal B)\le C_\phi\nu_Z(\mathcal B)
 \label{eq:representationcoverage}
\end{equation}
for every measurable representation event $\mathcal B$.
\end{assumption}

Histories merged at the same representation can still have different action errors.
Define
\begin{align}
 e(h)&=\norm{\pi(h)-\pi_{\mathrm E}(h)},\nonumber\\
 m_\rho(z)&=\E_\rho[e(h)^2\mid Z=z],\quad \rho\in\{\mu,\nu\},\nonumber\\
 D_\phi&=\E_{z\sim\mu_Z}|m_\mu(z)-m_\nu(z)|.
 \label{eq:conditionalerrordiscrepancy}
\end{align}
Assuming square-integrable action error and measurable conditional means, representation coverage gives
\begin{equation}
 \E_{\mu}e(h)^2
 \le C_\phi R_\nu(\pi)+D_\phi.
 \label{eq:representationrisktransfer-main}
\end{equation}
If Assumption~\ref{ass:actionsensitivity} also holds, the performance-difference identity gives
\begin{equation}
 F(\pi)\le F(\pi_{\mathrm E})
       +TL\sqrt{C_\phi R_\nu(\pi)+D_\phi}.
 \label{eq:representationcoveragebound}
\end{equation}

Appendix~\ref{app:representationcoverage} gives the proof.
Equation~\eqref{eq:representationcoveragebound} shows why improved representation coverage alone is insufficient.
A useful representation must reduce $C_\phi$ without increasing the conditional-error discrepancy $D_\phi$ enough to offset that gain.
